%!TEX program = xelatex
\documentclass[letterpaper,12pt]{article}
\usepackage{fontspec}\defaultfontfeatures{Ligatures=TeX}
\usepackage[letterpaper,vmargin={1in,1.5in},hmargin={1in,1in}]{geometry}
%-----------------------------------------------------------------------------%
\usepackage{settings}
%-----------------------------------------------------------------------------%
%%% Title %%%
\title{G6411 Recitation: Prerequisite Review Questions}
\author{Jesse Chieh Chen\\\href{mailto:cc5229@columbia.edu}{cc5229@columbia.edu}}
\date{\today}
%-----------------------------------------------------------------------------%

\begin{document}

\maketitle

\section{Probability and Statistics}

\begin{enumerate}
	\item(\textsc{Random Variable})
		Suppose you flip a fair coin that can land either heads ($H$) or tails ($T$).
		If it lands heads, the value $1$ is assigned as the outcome;
		if it lands tails, the value $-1$ is assigned.
		Let $X$ be the random variable representing the outcome.
		\begin{enumerate}
			\item
				What is the sample space of the experiment?
			\item
				Taking the event space to be the power set of the sample space,
				what are all the possible events?
			\item
				What is the probability of each event?
			\item
				What is a random variable?
				Define $X$ as a function on the sample space.
		\end{enumerate}

	\item(\textsc{Joint, Marginal, and Conditional Distributions})
		The joint probability mass function of $(X,Y)$ is given by
		\begin{center}
			\begin{tabular}{c|cc}
				& $Y=0$ & $Y=1$ \\
				\hline
				$X=0$ & $0.10$ & $0.20$ \\
				$X=1$ & $0.30$ & $0.40$
			\end{tabular}
		\end{center}
		\begin{enumerate}
			\item
				Verify that the table defines a valid joint distribution.
			\item
				Find the marginal distributions of $X$ and $Y$.
			\item
				Find $\Pr(X=1\given Y=1)$ and $\Pr(Y=1\given X=0)$.
			\item
				Compute $\E(X)$, $\var(X)$, $\E(Y)$, and $\var(Y)$.
			\item
				Compute $\cov(X,Y)$.
				Are $X$ and $Y$ independent?
		\end{enumerate}
	\item(\textsc{Continuous Random Variables})
		Suppose that the random variable $U$ has probability density function
		\begin{align}
			f_U(u)=
			\begin{cases}
				2u, & 0\leq u\leq 1,\\
				0, & \text{otherwise}.
			\end{cases}
		\end{align}
		\begin{enumerate}
			\item
				Verify that $f_U$ is a valid \gls{pdf}.
			\item
				Derive the \gls{cdf} $F_U(u)$, including its values outside $[0,1]$.
			\item
				Compute $\Pr(1/4<U\leq 3/4)$ using both $f_U$ and $F_U$.
			\item
				What is $\Pr(U=1/2)$?
			\item
				Find the median $m$ satisfying $F_U(m)=1/2$.
			\item
				Compute $\E(U)$, $\E(U^2)$, and $\var(U)$.
		\end{enumerate}
		Now suppose that the random variable $V$ has \gls{cdf}
		\begin{align}
			F_V(v)=
			\begin{cases}
				1-e^{-v}, & v\geq 0, \\
				0, & \text{otherwise.}
			\end{cases}
		\end{align}
		\begin{enumerate}[resume]
			\item
				Derive the \gls{pdf} $f_V(v)$.
			\item
				Compute $\Pr(V>2)$.
			\item
				Find the $90$th percentile of $V$.
			\item
				Compute $\E(V)$, $\E(V^2)$, and $\var(V)$.
		\end{enumerate}
\end{enumerate}

\section{Linear Algebra and Matrix Algebra}

Let
\begin{align}
	A=
	\begin{pmatrix}
		1 & 0 \\
		1 & 1 \\
		1 & 2
	\end{pmatrix}.
\end{align}

\begin{enumerate}
	\item
		What are the dimensions of $A$ and $A\T$? $A\T$ denotes the transpose of $A$.
	\item
		What is the rank of $A$?
	\item
		Compute $A\T A$.
	\item
		Is $A\T A$ invertible?
		If so, compute $(A\T A)\inv$.
	\item
		What property of $A$ guarantees that $A\T A$ is invertible?
\end{enumerate}

\section{Calculus: Taylor Expansions}

Let $g(x)=\log(x)$ for $x>0$.

\begin{enumerate}
	\item
		Compute $g'(x)$, $g''(x)$, and $g'''(x)$.
	\item
		State precisely what $o(h^k)$ and $O(h^k)$ mean as $h\to0$ for any fixed positive integer $k$.
		Also explain intuitively what these notations mean.
	\item
		Write the first-order and second-order Taylor expansions of $g(x)$ around $x=1$.
		Precisely, define $h=x-1$ and write the Taylor expansions of $g(1+h)$.
		For each expansion, express the remainder using both little-$o$ notation and the sharp big-$O$ order.
	\item
		Use both expansions to approximate $\log(1.1)$.
	\item
		Which approximation is more accurate?
		Relate the approximation errors to the remainder orders.
\end{enumerate}

\section{Calculus: Optimization}

Let $q:\reals^2\to\reals$ be defined by
\begin{align}
	q(x,y)=4x+2y-x^2-xy-y^2.
\end{align}

\begin{enumerate}
	\item(\textsc{Optimality Conditions})
		\begin{enumerate}
			\item
				State the \emph{first-order necessary condition} for an interior local maximum of a differentiable function.
			\item
				State the \emph{second-order sufficient condition} for an interior point to be a local maximum of a twice-differentiable function.
		\end{enumerate}
	\item
		Compute the gradient of $q$ and solve the first-order conditions.
	\item
		Compute the Hessian of $q$ and determine whether it is positive definite, negative definite, or indefinite.
	\item
		Use the second-order condition to classify the stationary point.
		Is it also the unique global maximizer?
\end{enumerate}

\end{document}
